\documentclass[../../main.tex]{subfiles}
\begin{document}

\begin{center}
  {\fontsize{48pt}{58pt}\selectfont \textbf{1986}\normalsize\textbf{年度}}
\end{center}

\begin{table}[htb]
  \centering
  \caption{出題テーマと難易度}
  \begin{tabular}{|c|p{9cm}|c|}
    \hline
    \textbf{問題番号} & \textbf{テーマ} & \textbf{難易度} \\
    \hline
    第1問 & 昇順に並ぶ実数の重み付き和 & \\
    \hline
    第2問 & 2つの正三角形を結ぶ線分が作る立体 & \\
    \hline
    第3問 & 正射影によって定まる一次変換 & \\
    \hline
    第4問 & 二つの三角形の中点と外接円 & \\
    \hline
    第5問 & サイコロの最大値と最小値の比 & \\
    \hline
    第6問 & 二つの円と四辺形の領域の面積 & \\
    \hline
  \end{tabular}
\end{table}

\begin{table}[htb]
  \centering
  \caption{解答}
  \renewcommand{\arraystretch}{1.4}
  \begin{tabular}{|c|c|p{9cm}|}
    \hline
    \textbf{問題番号} & \textbf{小問} & \textbf{答え} \\
    \hline
    第1問 & -- & $\sum_{i=1}^{n}ia_i>0$（証明問題） \\
    \hline
    第2問 & (1) & 切り口は一辺$\sqrt{3a^2-3a+1}$の正三角形 \\
    \cline{2-3}
    & (2) & $\dfrac{\sqrt{3}}{8}$ \\
    \hline
    第3問 & -- &  \\
    \hline
    第4問 & -- &  \\
    \hline
    第5問 & (1) &  \\
    \cline{2-3}
    & (2) &  \\
    \hline
    第6問 & -- & 最大面積は$1-\dfrac{\pi}{4}$ \\
    \hline
  \end{tabular}
\end{table}

\paragraph{試験形態}
6問．試験時間は要確認．

\end{document}
